\SourcePage{96}
\SGAchapter{XIII}{固有态射的基变换定理：证明的结束}{Théorème de changement de base pour un morphisme propre : fin de la démonstration}
\begin{center}M. Artin\end{center}
\SGAsection{1}{射影且平坦的情形}
回顾第 XII 讲 5.9 bis 在这里所考虑的情形中的陈述：
\begin{SGAresult}{命题}{1.1}
设 $S$ 是一个\Term{henselian}{亨泽尔}\Term{noetherian-ring}{诺特环}的谱，$f:X\to S$ 是\Term{projective}{射影}且\Term{flat}{平坦}的\Term{morphism}{态射}。设 $X_\circ$ 是 $X/S$ 的\Term{closed-fiber}{闭纤维}。\Term{restriction}{限制}\Term{functor}{函子}
\[
\operatorname{Et}(X)\longrightarrow\operatorname{Et}(X_\circ)
\]
是\Term{category-equivalence}{范畴等价}。
\end{SGAresult}
如前所述，由第 XII 讲 6.5 (i) 可知此\Term{arrow}{箭头}是\Term{fully-faithful}{全忠实}的。只须证明 $X_\circ$ 的每个\Term{etale-covering}{平展覆盖} $Y_\circ$ 都由 $X$ 的一个\Term{etale-covering}{平展覆盖} $Y$ 诱导。
\begin{SGAresult}{引理}{1.2}
设 $S$ \Term{locally-noetherian}{局部诺特}，$f:X\to S$ \Term{projective}{射影}且\Term{flat}{平坦}，$Y\subset X$ 是一个子\Term{scheme}{概形}。\Term{final-functor}{终函子}的\Term{subfunctor}{子函子}
\[
\varepsilon=\prod_{X/S}Y/X:(\mathrm{Sch})/S\longrightarrow(\mathrm{Ens}),
\]
“表示条件 $Y=X$”，即对于 $S'\to S$ 有
\[
\varepsilon(S')=\begin{cases}
\varnothing&\text{若 }Y'\ne X',\\
\{\varnothing\}&\text{若 }Y'=X',
\end{cases}
\]
其中撇号表示\Term{base-change}{基变换} $S'\to S$ 的结果。此\Term{functor}{函子}可由 $S$ 的一个子\Term{prescheme}{预概形} $Z$ 表示。若 $Y$ 是开的（相应地，闭的），则 $Z$ 也如此。
\end{SGAresult}
\SourcePage{97}
\begin{proof}
设 $Y$ 在 $X$ 的\Term{open-subset}{开子集} $U$ 中闭，$C$ 是 $U$ 的闭补集。由于 $f$ \Term{proper}{固有}，$C$ 在 $S$ 中的像是闭集 $D$；显然条件 $U=X$ 由 $S$ 的\Term{open-subset}{开子集} $S-D$ 表示。因此可以假定已经有 $U=X$，即 $Y$ 是闭子\Term{scheme}{概形}。考虑 $\mathcal O_X$-模的\Term{epimorphism}{满态射} $\mathcal O_X\xrightarrow{\alpha}\mathcal O_Y$；$Y=X$ \Term{equivalence}{等价}于存在\Term{morphism}{态射} $\mathcal O_Y\xrightarrow{\beta}\mathcal O_X$ 使复合 $\beta\alpha$ 为恒等\Term{morphism}{态射}，即存在 $\mathcal H om(\mathcal O_Y,\mathcal O_X)$ 的\Term{section}{截面}，\Term{lifting}{提升} $\mathcal H om(\mathcal O_X,\mathcal O_X)$ 的恒等\Term{section}{截面}。由 EGA III 7.7.8，把 $S'/S$ 映到 $\mathcal H om(\mathcal O_Y,\mathcal O_X)$（相应地，$\mathcal H om(\mathcal O_X,\mathcal O_X)$）的\Term{section}{截面}\Term{set}{集合}的\Term{functor}{函子}，由 $S$ 上的\Term{vector-bundle}{向量丛} $V_1$（相应地，$V_2$）表示。设 $\varphi:V_1\to V_2$ 是由 $\mathcal O_X\to\mathcal O_Y$ 诱导的\Term{morphism}{态射} $u\mapsto u\alpha$，$s:S\to V_2$ 是对应于恒等\Term{morphism}{态射} $\mathcal O_X\to\mathcal O_X$ 的\Term{section}{截面}。\Term{section}{截面} $s$ 在 $\varphi$ 下的\Term{inverse-image}{逆像}可等同于 $S$ 的一个闭子\Term{scheme}{概形} $Z$；显然 $Z$ 表示\Term{functor}{函子} $\varepsilon$，引理得证。
\end{proof}
\begin{SGAresult}{引理}{1.3}
设 $S$ 是\Term{locally-noetherian}{局部诺特}\Term{scheme}{概形}，$f:X\to S$ 是\Term{projective}{射影}且\Term{flat}{平坦}的\Term{morphism}{态射}，$E$ 是 $X$ 上的\Term{locally-free}{局部自由}\Term{sheaf}{层}。设
\[
QE:(\mathrm{Sch})/S\longrightarrow(\mathrm{Ens})
\]
是把 $S'/S$ 映到 $E'=E\otimes_{\mathcal O_S}\mathcal O_{S'}$ 的商\Term{sheaf}{层} $A'$ 的\Term{set}{集合} $QE(S')$ 的\Term{functor}{函子}，其中 $A'$ 配备 $\mathcal O_{X'}$ 上的\Term{etale}{平展}\Term{algebraic-structure}{代数结构}。此\Term{functor}{函子}可由 $S$ 上一个\Term{locally-finite-type}{局部有限型}\Term{scheme}{概形}表示。
\end{SGAresult}
\begin{proof}
\Term{functor}{函子} $\operatorname{Quot}(E)$，其中 $\operatorname{Quot}(E)(S')$ 是 $E'$ 的在 $S'$ 上\Term{flat}{平坦}的商的\Term{set}{集合}，可由一个\Term{locally-finite-type}{局部有限型}\Term{scheme}{概形}表示（TDTE IV，第 221 讲），记为 $Z/S$。设 $M_u$ 是 $X_Z$ 上 $E\otimes_{\mathcal O_S}\mathcal O_Z$ 的“泛”商。$\operatorname{Quot}(E)$ 中由在 $X'$ 上\Term{flat}{平坦}的商所构成的\Term{subfunctor}{子函子}，由 $Z$ 的一个\Term{open-subset}{开子集} $U$ 表示。事实上，设 $M'$ 是一个商 $E'$，由 $S$-\Term{morphism}{态射} $S'\to Z$ 给出。由于 $M'$ 在 $S'$ 上\Term{flat}{平坦}，它在 $X'$ 上\Term{flat}{平坦}\Term{equivalence}{等价}于：对于每个 $s'\in S'$，模 $M'_{s'}$ 在\Term{fiber}{纤维} $X'_{s'}$ 上\Term{flat}{平坦}（EGA IV 5.9）；这又\Term{equivalence}{等价}于 $(M_u)_z$ 在\Term{fiber}{纤维} $Z_z$ 上\Term{flat}{平坦}，其中 $z$ 是 $s'$ 的像。由 EGA IV 11.1.5，满足 $(M_u)_z$ 在 $Z_z$ 上\Term{flat}{平坦}，即 $M_u$ 在 $Z_s$ 的各点处在 $Z$ 上\Term{flat}{平坦}的 $z\in Z$ 构成 $Z$ 的\Term{open-subset}{开集} $U$，所言得证。

现在只须表示相对\Term{functor}{函子} $QE/U[f]$。因此，以 $U$ 代替 $S$、以 $X\times_S U$ 代替 $X$，就归结为下一个引理。
\end{proof}
\begin{SGAresult}{引理}{1.4}
设 $S$ 是\Term{locally-noetherian}{局部诺特}\Term{scheme}{概形}，$f:X\to S$ 是\Term{projective}{射影}且\Term{flat}{平坦}的\Term{morphism}{态射}，$M$ 是 $\mathcal O_X$ 上的\Term{locally-free}{局部自由}\Term{sheaf}{层}。\Term{functor}{函子}
\[
\operatorname{Alg\text{é}t}:(\mathrm{Sch})/S\longrightarrow(\mathrm{Ens}),
\]
把 $S'/S$ 映到 $M'$ 上在 $\mathcal O_{X'}$ 上\Term{etale}{平展}的\Term{algebraic-structure}{代数结构}的\Term{set}{集合} $\operatorname{Alg\text{é}t}(S')$，可由 $S$ 上一个\Term{finite-type}{有限型}\Term{scheme}{概形}表示。
\end{SGAresult}
\SourcePage{98}
\begin{proof}
为简洁起见，说“\Term{representable-functor}{可表示}”时省略“由一个相对\Term{finite-type}{有限型}\Term{scheme}{概形}”几个字。所出现的\Term{vector-bundle}{向量丛}都由\Term{coherent}{凝聚}\Term{sheaf}{层}定义，因而都是\Term{finite-type}{有限型}的。$M'$ 上的\Term{bilinear}{双线性}合成律由\Term{locally-free}{局部自由}\Term{sheaf}{层} $\mathcal H om(M'\otimes M',M')$ 的一个\Term{section}{截面}给出。由于 $\mathcal H om(M'\otimes M',M')$ 是 $S$-\Term{flat}{平坦}的，由 EGA III 7.7.6，“\Term{bilinear}{双线性}合成律”\Term{functor}{函子}由 $S$ 上的一个\Term{vector-bundle}{向量丛} $Z$ 表示。须验证此合成律满足结合律、交换律以及存在\Term{unit-section}{单位截面}这几个附加条件，由 $Z$ 的一个子\Term{scheme}{概形}表示。考虑相对\Term{functor}{函子}，可以用 $Z$ 代替 $S$，从而假定 $M$ 上已给定合成律，并证明表示上述条件的 $S$ 的\Term{subfunctor}{子函子}可由 $S$ 的一个子\Term{scheme}{概形}表示。

对于结合律，条件是 $\mathcal H om(M\otimes M\otimes M,M)$ 的两个\Term{section}{截面}在\Term{base-change}{基变换} $S'\to S$ 后相等。设 $Y$ 是这两个\Term{section}{截面}在对应于 $\mathcal H om(M\times M\otimes M,M)$ 的\Term{vector-bundle}{向量丛}上的\Term{fiber-product}{纤维积}。\Term{scheme}{概形} $Y$ 可等同于 $X$ 的一个闭子\Term{scheme}{概形}；$S'\to S$ 后的结合律由条件 $Y'=X'$ 表示。由 1.2，它定义一个可由 $S$ 的闭子\Term{scheme}{概形}表示的 $S$ 的\Term{subfunctor}{子函子}。对于交换律，条件是 $\mathcal H om(M\otimes M,M)$ 的两个\Term{section}{截面}变得相等，由 1.2，这同样定义可由闭子\Term{scheme}{概形}表示的 $S$ 的\Term{subfunctor}{子函子}。最后，现在可以假定此合成律已满足结合律和交换律，考察存在\Term{unit-section}{单位截面}的条件；此\Term{section}{截面}若存在便唯一。由 EGA III 7.7.6，把 $S'/S$ 映到 $M'$ 的\Term{section}{截面}\Term{set}{集合}的\Term{functor}{函子}可由一个\Term{vector-bundle}{向量丛}表示。因此，只须表示此\Term{functor}{函子}中对应于\Term{unit-section}{单位截面}的\Term{subfunctor}{子函子}（注意刚才提到的唯一性）。照常相对化，可以假定 $M$ 的一个\Term{section}{截面} $\varepsilon$ 已经给定，并证明表示它变成\Term{unit-section}{单位截面}的 $S$ 的\Term{subfunctor}{子函子}\Term{representable-functor}{可表示}。$\varepsilon$ 为单位的条件是显然的\Term{morphism}{态射}的复合
\[
M'\longrightarrow M'\otimes M'\longrightarrow M'
\]
为恒等\Term{morphism}{态射}，即 $\mathcal H om(M,M)$ 的两个确定的\Term{section}{截面}变得相等。再用 1.2 即可。

最后，假定\Term{locally-free}{局部自由}模 $M$ 上已经给定交换含幺\Term{algebraic-structure}{代数结构}。\Term{etale}{平展}条件在 $X$ 中是开条件（SGA I 4.5），因而在 $S$ 上可由一个（开）子\Term{scheme}{概形}表示。这就证明了 1.4，进而证明了 1.3。
\end{proof}
\textbf{1.5。}现在设 $Z/S$ 是表示 1.3 中\Term{functor}{函子}的\Term{object}{对象}，考察 $Z$ 在一点 $z$ 处在 $S$ 上\Term{smooth}{光滑}的条件。将 SGA I 第 III 讲 3.1 翻译成\Term{functor}{函子} $QE$ 的语言，得到如下条件。

设 $B$ 是\Term{artinian-local-ring}{阿廷局部环}，$\mathfrak m=\operatorname{rad}B$，$J$ 是满足 $J\mathfrak m=0$ 的 $B$ 的\Term{ideal}{理想}，$S'=\operatorname{Spec}B$，$S''=\operatorname{Spec}(B/J)$，并给定\Term{morphism}{态射} $S'\to S$。对\Term{base-change}{基变换} $S'\to S$ 与 $S''\to S$ 的结果，采用显然的记号。
\SourcePage{99}

$(*)$ 对于 $E''$ 的每个配有 $\mathcal O_{X''}$ 上\Term{etale}{平展}\Term{algebraic-structure}{代数结构}、并对应于像为 $z$ 的 $S$-\Term{morphism}{态射} $S''\to Z$ 的商 $A''$，存在 $E'$ 的一个配有 $\mathcal O_{X'}$ 上\Term{etale}{平展}\Term{algebraic-structure}{代数结构}的商 $A'$，诱\Term{derive}{导出}它。

由于 $X'$ 和 $X''$ 具有相同的底\Term{space}{空间}，诱导 $A''$ 的\Term{algebra}{代数} $A'$ 已经由此唯一确定到唯一\Term{isomorphism}{同构}（SGA I 8.3）。因此上述条件\Term{representable}{可表}为：

$(**)$ 设 $A'$ 是 $\mathcal O_{X'}$ 上的\Term{etale}{平展}\Term{algebra}{代数}，$A''$ 是 $A'$ 在 $X''$ 上诱导的\Term{algebra}{代数}。每个对应于像为 $z$ 的\Term{morphism}{态射} $S''\to Z$ 的满\Term{homomorphism}{同态}
\[
E''\xrightarrow{\varphi''}A''
\]
都可以\Term{lifting}{提升}为一个\Term{homomorphism}{同态}
\[
E'\xrightarrow{\varphi'}A',
\]
后者由 \Term{nakayama-lemma}{Nakayama 引理}必为\Term{surjectivity}{满射}。

由\Term{exact-sequence}{正合列}
\[
0\longrightarrow J\longrightarrow B\longrightarrow B/J\longrightarrow0
\]
得到 $X'$ 上的\Term{sheaf}{层}的\Term{exact-sequence}{正合列}
\[
0\longrightarrow J\otimes\mathcal H om(E',A')\longrightarrow\mathcal H om(E',A')\longrightarrow\mathcal H om(E'',A'')\longrightarrow0.
\]
\Term{lifting}{提升} $\varphi''$ 的障碍属于 $H^1(X',J\otimes\mathcal H om(E',A'))$。设 $S_\circ=\operatorname{Spec}(B/\mathfrak m)$。由于 $J\mathfrak m=0$，有
\[
J\otimes\mathcal H om(E',A')\simeq J\otimes\mathcal H om(E_\circ,A_\circ),
\]
而且显然
\[
H^1(X',J\otimes\mathcal H om(E_\circ,A_\circ))\simeq J\otimes H^1(X_\circ,\mathcal H om(E_\circ,A_\circ)).
\]
考虑到\Term{cohomology}{上同调}与\Term{base-field-extension}{基域扩张}可交换，得到：
\begin{SGAresult}{推论}{1.6}
设 $Z/S$ 是表示 1.3 中\Term{functor}{函子}的\Term{scheme}{概形}，$X_Z=X\times_S Z$，$A_Z$ 是 $E_Z$ 的\Term{etale}{平展}\Term{algebra}{代数}商（“一般”商）。则 $Z/S$ 在满足
\[
H^1(X_\circ,\mathcal H om(E_\circ,A_\circ))=0
\]
的每个点 $z_\circ$ 处\Term{smooth}{光滑}；这里下标零表示\Term{restriction}{限制}到 $X_Z/Z$ 在 $z_\circ$ 处的\Term{fiber}{纤维} $X_\circ$。
\end{SGAresult}
现在可以完成定理 1.1 的证明。采用定理中的记号，设 $Y_\circ/X_\circ$ 是\Term{etale-covering}{平展覆盖}，由 $\mathcal O_{X_\circ}$ 上\Term{etale}{平展}的\Term{algebra}{代数}\Term{sheaf}{层} $A_\circ$ 给出。对于充分大的 $n$，$A_\circ(n)$ 由其\Term{section}{截面}生成（EGA III 2.2.1）。因此有\Term{surjectivity}{满射}
\[
\mathcal O_{X_\circ}^{N}\longrightarrow A_\circ(n)\longrightarrow0,
\]
从而有\Term{surjectivity}{满射}
\[
\mathcal O_{X_\circ}^{N}(-n)=E_\circ\longrightarrow A_\circ\longrightarrow0.
\]
\SourcePage{100}
取充分大的 $n$，还使得
\[
H^1(X_\circ,\mathcal H om(E_\circ,A_\circ))=0,\tag{+}
\]
由 EGA III 2.2.1 这是可能的，因为
\[
\mathcal H om(E_\circ,A_\circ)=[\mathcal H om(\mathcal O_{X_\circ}^{N},A_\circ)](n)=[A_\circ(n)]^N.
\]
置 $E=\mathcal O_X^N(-n)$，设 $Z/S$ 是表示 1.3 中\Term{functor}{函子}的\Term{scheme}{概形}。$E_\circ$ 的商 $A_\circ$ 对应于 $Z/S$ \Term{closed-fiber}{闭纤维}中的一点 $z_\circ\in Z$，它在 $k(s)$ 上有理；由 1.5 与 $(+)$，$Z/S$ 在 $z_\circ$ 处\Term{smooth}{光滑}。由于 $S$ 是\Term{henselian}{亨泽尔}环的谱，且 $Z/S$ 在\Term{closed-fiber}{闭纤维}的点 $z_\circ$ 处\Term{smooth}{光滑}，可以找到 $Z/S$ 的一个经过 $z_\circ$ 的\Term{section}{截面}；例如，这由\Term{smooth}{光滑}性的原始定义（SGA II 1.1）立即得到。因此 $A_\circ$ 由 $X$ 上的一个\Term{etale}{平展}\Term{algebra}{代数} $A$（$E$ 的商）诱导，1.1 的证明完成。

\SGAsection{2}{相对维数不超过一的情形}
所要证明的是：
\begin{SGAresult}{命题}{2.1}
采用第 XII 讲 5.9 bis 的记号，若 $S$ \Term{noetherian}{诺特}且\Term{strictly-local}{严格局部}，$f$ \Term{projective}{射影}且\Term{dimension}{维数} $\leq1$，则\Term{functor}{函子} $\operatorname{Et}(X)\to\operatorname{Et}(X_\circ)$ 是\Term{category-equivalence}{范畴等价}。
\end{SGAresult}
根据第 XII 讲 8.5，这将完成第 XII 讲 5.1 (ii) 的证明；后者同时推广了 1.1 和 2.1。
\begin{SGAresult}{引理}{2.2}
设 $S$ 是\Term{noetherian-local-ring}{诺特局部环} $A$ 的谱，$X/S$ \Term{projective}{射影}，\Term{relative-dimension}{相对维数} $\leq n$。存在\Term{relative-dimension}{相对维数} $\leq n$ 的\Term{projective}{射影}且\Term{flat}{平坦}的 $S$-\Term{scheme}{概形} $Z/S$，以及 $S$ 上的一个\Term{closed-immersion}{闭浸入} $X\hookrightarrow Z$。
\end{SGAresult}
\begin{proof}
设 $X\to\mathbf P_S^N$ 是 $S$ 上的一个\Term{projective}{射影}\Term{embedding}{嵌入}，$J\subset A[x_\circ,\ldots,x_N]$ 是定义 $X$ 的\Term{homogeneous}{齐次}\Term{ideal}{理想}。设 $k=A/\operatorname{rad}A$，$S_\circ=\operatorname{Spec}k$，$X_\circ=X\otimes_S k$。则 $X_\circ$ 由\Term{homogeneous}{齐次}\Term{ideal}{理想} $J_\circ=\operatorname{Im}(J\to k[X_\circ,\ldots,X_N])$ 定义。由于 $\dim X_\circ\leq n$，众所周知，存在 $J_\circ$ 的元素 $f_1^\circ,\ldots,f_{N-n}^\circ$，定义 $\mathbf P_k^N$ 的一个 $n$ 维子\Term{scheme}{概形} $Z_\circ$。取 $J$ 中\Term{lifting}{提升} $f_i^\circ$ 的元素 $f_i$（$0\leq i\leq N-n$），并设 $Z$ 是由 $f_1,\ldots,f_\nu$ 定义的 $\mathbf P_S^N$ 的子\Term{scheme}{概形}。由 EGA IV 11.3，$Z$ 在 $S$ 上\Term{flat}{平坦}；另一方面，由 EGA IV 13.1.5，其\Term{relative-dimension}{相对维数}为 $\dim Z_\circ=n$。2.2 得证。
\end{proof}
\begin{proof}[2.1 的证明]
设 $X/S$ \Term{projective}{射影}、\Term{relative-dimension}{相对维数} $\leq1$，$S$ \Term{noetherian}{诺特}且\Term{strictly-local}{严格局部}。可以假定 $X$（从而 $X_\circ$）\Term{connected}{连通}。取 $S$ 上一个\Term{closed-immersion}{闭浸入} $X\to Z$，其中 $Z$ \Term{projective}{射影}、\Term{flat}{平坦}且\Term{dimension}{维数} $\leq1$，见 2.2。设 $s_\circ$ 是 $S$ 的\Term{closed-point}{闭点}。有相应的图表
\SourcePage{101}
\[
\begin{tikzcd}
\operatorname{Et}(Z)\arrow[r]\arrow[d,"a"']&\operatorname{Et}(X)\arrow[d,"b"]\\
\operatorname{Et}(Z_\circ)\arrow[r,"c"']&\operatorname{Et}(X_\circ).
\end{tikzcd}
\]
须证明 $b$ 是\Term{equivalence}{等价}。由第 XII 讲 5.8，$b$ \Term{fully-faithful}{全忠实}，只须证明它\Term{essentially-surjective}{本质满射}。由 1.1，$a$ \Term{essentially-surjective}{本质满射}，所以只须证明 $c$ \Term{essentially-surjective}{本质满射}。略去下标 $\circ$，即归结为下述命题。
\end{proof}
\begin{SGAresult}{命题}{2.3}
设 $k$ 是\Term{separably-closed-field}{可分闭域}，$Z$ 是\Term{locally-finite-type}{局部有限型}、\Term{dimension}{维数} $\leq1$ 的 $k$-\Term{scheme}{概形}，$X\hookrightarrow Z$ 是\Term{connected}{连通}闭子\Term{scheme}{概形}。则 $X$ 的每个\Term{etale-covering}{平展覆盖}都由 $Z$ 的一个\Term{etale-covering}{平展覆盖}诱导。
\end{SGAresult}
\begin{proof}
可以假定 $Z$ 和 $X$ 都为一维；$X$ 零维的情形由对 $k$ 的假设是平凡的。于是 $Z=X\cup Y$，其中 $Y$ 是 $Z-X$ 的\Term{closure}{闭包}所给出的闭子\Term{scheme}{概形}。设 $V=X\cap Y$ 是交\Term{scheme}{概形}，它是一个\Term{discrete}{离散}\Term{scheme}{概形}，即\Term{local-artinian}{局部阿廷}\Term{scheme}{概形}。\Term{morphism}{态射} $X\amalg Y\to Z$ \Term{finite}{有限}且\Term{surjectivity}{满射}，因而是\Term{etale-covering}{平展覆盖}\Term{category}{范畴}的严格\Term{descent}{下降}\Term{morphism}{态射}（SGA IX 4.7）。显然
\[
(X\amalg Y)\times_Z(X\amalg Y)=X\amalg Y\amalg V^{(1)}\amalg V^{(2)}\qquad(V^{(i)}=V),\tag{*}
\]
其中 $\operatorname{pr}_i$ 把 $V^{(i)}$ 映入 $X\amalg Y$ 的分支 $X$，而在 $j\ne i$ 时把 $V^{(j)}$ 映入分支 $Y$。

设 $X'/X$ 是次数 $n$ 的\Term{etale-covering}{平展覆盖}，$Y''/Y$ 是次数 $n$ 的\Term{covering}{覆盖}（例如次数 $n$ 的\Term{trivial-covering}{平凡覆盖}），从而 $X'\amalg Y''$ 是 $X\amalg Y$ 的次数 $n$ 的\Term{covering}{覆盖}。只须为这个\Term{covering}{覆盖}及\Term{morphism}{态射} $X\amalg Y\to Z$ 找到\Term{descent-data}{下降资料}。由 $(*)$ 可见，此\Term{morphism}{态射}的一个\Term{descent-data}{下降资料}恰好相当于一个\Term{isomorphism}{同构}
\[
X'\times_X V=V'\xrightarrow{\sim}V''=Y''\times_Y V.
\]
由于 $k$ \Term{separably-closed}{可分闭}，$k(z_i)$ 也如此，因而 $V$ 的每个\Term{etale-covering}{平展覆盖}都\Term{completely-split}{完全分裂}，故 $V'\simeq V''$，结论得证。
\end{proof}
\textbf{注记 2.4。}可以用直接论证避免引用 SGA IX 中这个精细的结果和\Term{descent}{下降}理论：证明 $Z$ 的\Term{etale-covering}{平展覆盖}\Term{category-equivalence}{范畴等价}于三元组 $(X',Y',\theta)$ 的\Term{category}{范畴}，其中 $X'$（相应地，$Y'$）是 $X$（相应地，$Y$）的\Term{etale-covering}{平展覆盖}，$\theta$ 是 $V$-\Term{isomorphism}{同构} $X'\times_X V\simeq Y'\times_Y V$。对于任意由两个闭子\Term{scheme}{概形} $X,Y$ 之并构成的\Term{scheme}{概形} $Z$，此结果均成立；可以直接利用确定 $X$ 上\Term{quasi-coherent}{拟凝聚}模的类似结果（当 $X=V(J)$，$Y=V(K)$，$J\cap K=0$ 时）证明。

\SourcePage{102}
\SGAsection{3}{关于 Picard 群的一个辅助结果}
\footnote{此结果也以稍广的\Term{form}{形式}见于 EGA IV 21.9.11、21.9.12；特别地，在 3.2 中只须 $f$ \Term{separated}{分离}，无须\Term{proper}{固有}。}
\begin{SGAresult}{命题}{3.1}
设 $X$ 是\Term{noetherian}{诺特}\Term{scheme}{概形}，$X_\circ$ 是 $X$ 的闭子\Term{scheme}{概形}，满足：a) $\dim X_\circ\leq1$；b) 对 $X$ 的每个\Term{nonempty}{非空}\Term{connected}{连通}\Term{closed-subset}{闭子集} $Y$，$Y_\circ=Y\cap X_\circ$ \Term{connected}{连通}且\Term{nonempty}{非空}。

(i) 设 $Z_\circ$ 是 $\overline{\{x\}}\cap X_\circ$ 的并，其中 $x\in\operatorname{Ass}\mathcal O_X$ 且 $\overline{\{x\}}\cap X_\circ$ 只含一个点。则对 $X_\circ$ 上\Term{support}{支撑}不交 $Z_\circ$ 的每个\Term{cartier-divisor}{Cartier 除子} $D_\circ$，存在 $X$ 上诱导 $D_\circ$ 的\Term{cartier-divisor}{Cartier 除子} $D$（且当 $D_\circ\geq0$ 时，可取 $D\geq0$）。

(ii) 假定存在 $X_\circ$ 的一个\Term{affine}{仿射}\Term{open-subset}{开子集} $U_\circ$，包含\Term{finite}{有限}集 $\operatorname{Ass}\mathcal O_{X_\circ}\cup Z_\circ$（当 $X$ \Term{noetherian}{诺特}且 $X_\circ$ 具有充足\Term{invertible}{可逆}\Term{sheaf}{层}时，此条件自动满足）。则\Term{homomorphism}{同态} $\operatorname{Pic}(X)\to\operatorname{Pic}(X_\circ)$ \Term{surjectivity}{满射}。
\end{SGAresult}
\begin{proof}
(i) 由于 $\dim X_\circ\leq1$，立即可见 $X_\circ$ 上的每个\Term{cartier-divisor}{Cartier 除子} $D_\circ$ 的\Term{support}{支撑}都是\Term{discrete}{离散}的，因而它是两个正\Term{cartier-divisor}{Cartier 除子}之差，二者的\Term{support}{支撑}包含于 $D_\circ$ 的\Term{support}{支撑}。因此 (i) 归结为 $D_\circ$ 为正的情形，即对于每个 $x\in\operatorname{Supp}D_\circ$，$D_\circ$ 在 $x$ 处由一个\Term{non-zero-divisor}{非零因子} $f_x\in\operatorname{rad}(\mathcal O_{X_\circ,x})$ 定义。$f_x$ 由 $\mathcal O_X$ 在 $x$ 的某个\Term{open-neighborhood}{开邻域} $U_x$ 上的一个\Term{section}{截面} $g_x$ 诱导。我断言，当 $U_x$ 取得充分小时，$g_x$ 是 $\mathcal O_{U_x}$ 中的\Term{non-zero-divisor}{非零因子}，即 $V(g_x)\cap\operatorname{Ass}\mathcal O_{U_x}=\varnothing$。事实上，由 a)、b)，对于每个 $z\in X$，特别是 $z\in\operatorname{Ass}\mathcal O_X$，$\overline{\{z\}}\cap X_\circ$ 是 $X_\circ$ 的一个\Term{connected}{连通}闭集，因而或者只含一个\Term{closed-point}{闭点}，或者是 $X_\circ$ 的若干一维\Term{irreducible-component}{不可约分支}的\Term{connected}{连通}并。由于 $x\notin Z_\circ$，在第一种情形有 $x\notin\overline{\{z\}}$。对第二种类型的 $z\in\operatorname{Ass}\mathcal O_X$，不可能同时有 $x\in\overline{\{z\}}$ 和 $z\in V(g_x)$，否则 $g_x$、从而 $f_x$ 将在 $\overline{\{z\}}\cap\operatorname{Supp}X_\circ$ 的某个\Term{irreducible-component}{不可约分支}上为零，该分支也是 $X_\circ$ 的\Term{irreducible-component}{不可约分支}，与 $f_x$ 在 $\mathcal O_{X_\circ,x}$ 中是\Term{non-zero-divisor}{非零因子}矛盾。因此缩小 $U_x$ 后，可假定 $V(g_x)\cap\operatorname{Ass}\mathcal O_{U_x}=\varnothing$，即 $g_x$ 为\Term{non-zero-divisor}{非零因子}。设 $Y_x$ 是 $V(g_x)$ 在 $X$ 中的\Term{closure}{闭包}，则 $Y_x\cap X_\circ\cap U_x\subset V(g_x)\cap X_\circ$ 以 $x$ 为\Term{isolated-point}{孤立点}，故 $Y_x\cap X_\circ$ 也如此。由命题的条件 b)，$Y_x$ 有一个\Term{connected-component}{连通分支} $Y'_x$，满足 $Y'_x\cap X_\circ=\{x\}$。而且 $Y'_x\subset U_x$，因为 $T_x=Y'_x-Y'_x\cap U_x$ 是 $X$ 的一个闭集，满足 $T_x\cap X_\circ=\varnothing$，由条件 b) 得 $T_x=\varnothing$。现在取 $X$ 上的\Term{cartier-divisor}{Cartier 除子} $D'_x$，使它在 $U_x$ 上诱导 $Y'_x$、在 $X-Y'_x$ 上诱导零，并置 $D=\sum_xD'_x$；它是 $X$ 上的一个正\Term{cartier-divisor}{Cartier 除子}，显然在 $X_\circ$ 上诱导 $D_\circ$。

(ii) 结论由 (i) 和以下事实得到：$X_\circ$ 上每个\Term{invertible-module}{可逆模} $L_\circ$ 都\Term{isomorphism}{同构}于一个\Term{support}{支撑}不交 $Z_\circ$ 的\Term{cartier-divisor}{Cartier 除子} $D_\circ$ 所定义的模。后一个事实事实上是说，$L_\circ$ 在包含 $\operatorname{Ass}\mathcal O_{X_\circ}\cup Z_\circ$ 的某个\Term{open-subset}{开集} $V_\circ$ 上有一个\Term{section}{截面}，且此\Term{section}{截面}在这个\Term{set}{集合}的每一点都不为零。
\SourcePage{103}
只须以\Term{affine-open}{仿射开集} $U_\circ$ 代替 $X_\circ$ 来证明此事，而在这种情形这是熟知的，来自半\Term{local-scheme}{局部}\Term{scheme}{概形}上的\Term{invertible-module}{可逆模}\Term{isomorphism}{同构}于结构\Term{sheaf}{层}这一事实。最后，众所周知，若 $X_\circ$ 具有充足\Term{sheaf}{层}，则 $X_\circ$ 的每个\Term{finite}{有限}\Term{subset}{子集}都包含于一个\Term{affine}{仿射}\Term{open-subset}{开子集}中。这就证明了 (ii)。
\end{proof}
以下推论推广第 XII 讲 8.6，并完成第 XII 讲 5.1 (iii) 的证明。
\begin{SGAresult}{推论}{3.2}
设 $S$ 是\Term{noetherian}{诺特}\Term{henselian}{亨泽尔}环的谱，$f:X\to S$ 是\Term{proper}{固有}且\Term{relative-dimension}{相对维数} $\leq1$ 的\Term{morphism}{态射}。则对于 $X$ 的每个与 $X/S$ 的\Term{closed-fiber}{闭纤维} $X_\circ$ 具有相同底\Term{space}{空间}的闭子\Term{scheme}{概形} $X'_\circ$，\Term{canonical}{典范}映射 $\operatorname{Pic}(X)\to\operatorname{Pic}(X_\circ)$ \Term{surjectivity}{满射}。
\end{SGAresult}
事实上，把命题用于 $(X,X'_\circ)$；条件 a) 由假设满足，条件 b) 由第 XII 讲 5.7 满足。只须验证 (ii) 中关于 $U_\circ$ 存在的假设；这是因为 $X'_\circ$ 必为\Term{projective}{射影}的（正如域上的每条\Term{proper}{固有}曲线一样）。

\textbf{注记。}不愿承认或重新证明域上的\Term{proper}{固有}曲线是\Term{projective}{射影}的读者，可以只在已经假定 $f$ \Term{projective}{射影}的情形使用此推论，这已经足够用于这里所需的应用。另一方面，若承认域上\Term{proper}{固有}\Term{algebraic-curve}{代数曲线}的\Term{projective}{射影}性，则在 3.2 的条件下，$X$ 必在 $S$ 上\Term{projective}{射影}：取一个充足的 $\xi_\circ\in\operatorname{Pic}(X_\circ)$，将它\Term{lifting}{提升}为 $\xi\in\operatorname{Pic}(X)$，再用 EGA III 4.7.1 即可。

\SGAsection{}{参考文献}
\begin{enumerate}
\item[{[1]}] A. Grothendieck，\emph{\Term{algebra}{代数}几何中的\Term{descent}{下降}技术与\Term{existence-theorem}{存在定理}}（Technique de descente et théorèmes d'existence en Géométrie Algébrique），载于\emph{\Term{algebra}{代数}几何基础}（Fondements de la Géométrie Algébrique，Bourbaki 讨论班 1957--1962 年摘录），Secrétariat mathématique，11, rue Pierre Curie。
\end{enumerate}

\SourcePage{104}
\SGAchapter{XIV}{固有态射的有限性定理；仿射代数概形的上同调维数}{Théorème de finitude pour un morphisme propre ; dimension cohomologique des schémas algébriques affines}
\begin{center}M. Artin\end{center}
本讲包含两个重要定理，它们都利用\Term{proper}{固有}\Term{morphism}{态射}的\Term{base-change-theorem}{基变换定理}来证明。
\SGAsection{1}{固有态射的有限性定理}
\begin{SGAresult}{定理}{1.1}
设 $f:X\to S$ 是\Term{proper}{固有}且\Term{finite-presentation}{有限表示}的\Term{morphism}{态射}。设 $A$ 是\Term{left-noetherian-ring}{左诺特环}，并对某个 $n\in\mathbf N$ 为 $\mathbf Z/n$-\Term{algebra}{代数}。设 $F$ 是 $X$ 上的\Term{constructible}{可构造}\Term{set}{集合}\Term{sheaf}{层}（相应地，群\Term{sheaf}{层}，$A$-模\Term{sheaf}{层}）。则\Term{sheaf}{层} $f_*F$（相应地，$R^1f_*F$，对每个 $q\geq0$ 的 $R^qf_*F$）也\Term{constructible}{可构造}。
\end{SGAresult}
特别地，有：
\begin{SGAresult}{推论}{1.2}
设 $X$ 是在\Term{separably-closed-field}{可分闭域} $k$ 的谱上\Term{proper}{固有}的\Term{scheme}{概形}，$F$ 是 $X$ 上的\Term{constructible}{可构造}\Term{torsion-sheaf}{挠阿贝尔层}。则对每个 $q\geq0$，群 $H^q(X,F)$ 都\Term{finite}{有限}。
\end{SGAresult}
\SourcePage{105}
\textbf{注记 1.3。}若去掉 $f$ \Term{proper}{固有}的条件，则该定理对阿贝尔\Term{sheaf}{层}不成立：取 $F=\mathbf Z/p$，$p=\operatorname{char}k$，$X=\mathbf E_k^1$ 为\Term{affine-space}{仿射空间}，$f$ 为 $X$ 到 $\mathbf P_k^1$ 的\Term{inclusion-morphism}{包含态射}，或\Term{structural-morphism}{结构态射} $X\to\operatorname{Spec}k$，即可看出。不过猜想当 $F$ 的阶与 $S$ 的\Term{residue-characteristic}{剩余特征}互素、$S$ “\Term{excellent}{优良}”（EGA IV 7.8）时，定理仍成立。对于\Term{equal-characteristic}{等特征}的\Term{excellent}{优良}\Term{scheme}{概形} $S$，若有\Term{resolution-singularities}{奇点消解}定理（例如在\Term{characteristic-zero}{特征零}时），便可证明此事（第 XIX 讲 5.1）。
\begin{proof}[1.1 的证明]
断言在 $S$ 上\Term{local-scheme}{局部}，因此可以取 $S$ \Term{affine}{仿射}。于是 $S$ 是 $\mathbf Z$ 上\Term{finite-type}{有限型}环的谱的\Term{limit}{极限}；定理的数据可以通过\Term{base-change}{基变换} $S\to S_\circ$，由一个\Term{proper}{固有}\Term{morphism}{态射} $f_\circ:X_\circ\to S_\circ$ 以及 $X_\circ$ 上的\Term{constructible}{可构造}\Term{sheaf}{层} $F_\circ$ 得到，其中 $S_\circ$ 在 $\mathbf Z$ 上\Term{finite-type}{有限型}，见 EGA IV 8 及第 IX 讲 2.7.4。由于\Term{constructible}{可构造}\Term{sheaf}{层}的\Term{inverse-image}{逆像}仍\Term{constructible}{可构造}（第 IX 讲 2.4 (iii)），把第 XII 讲 5.1 用于\Term{base-change-morphism}{基变换态射} $S_\circ\leftarrow S$，便归结为证明 $f_{\circ*}F_\circ$（相应地，$R^qf_{\circ*}F_\circ$）\Term{constructible}{可构造}，从而得到下一个引理。
\end{proof}
\begin{SGAresult}{引理}{1.4}
只须在 $S$ 为 $\operatorname{Spec}\mathbf Z$ 上\Term{finite-type}{有限型}的情形验证 1.1。
\end{SGAresult}
\begin{SGAresult}{引理}{1.5}
1.1 关于\Term{set}{集合}的断言成立，即若 $F$ 是\Term{constructible}{可构造}\Term{set}{集合}\Term{sheaf}{层}，则 $f_*F$ 也是。
\end{SGAresult}
\begin{proof}
可以假定 $S$ \Term{noetherian}{诺特}（1.4）。取一个\Term{injectivity}{单射} $F\to G=\prod_i\pi_{i*}(C_i)$，其中 $\pi_i:X_i\to X$ \Term{finite}{有限}，$C_i$ 在 $X_i$ 上常值且\Term{constructible}{可构造}（第 IX 讲 2.14）。有 $f_*F\subset f_*G$，因此由第 IX 讲 2.9 (iii)，只须处理 $F=G$ 的情形。但 $f_*G=\prod_i(f\pi_i)_*G_i$，且 $S$ 上\Term{constructible}{可构造}\Term{sheaf}{层}的\Term{finite-product}{有限积}仍\Term{constructible}{可构造}（第 IX 讲 2.14）。以 $X_i$ 代替 $X$，便归结为 $F=D_X$ 是取值 $D$ 的常值\Term{constructible}{可构造}\Term{sheaf}{层}的情形。由第 XII 讲 5.1 (i)，\Term{fiber}{纤维} $(f_*F)_{\bar s}$ \Term{isomorphism}{同构}于 $D^{C(\bar s)}$，其中 $C(\bar s)=\Pi_0(X_{\bar s})$，故此\Term{fiber}{纤维}\Term{finite}{有限}。由第 IX 讲 2.13 (iii)，只须证明函数“$f_*F$ 的\Term{fiber}{纤维}的元素个数”在 $S$ 上\Term{constructible}{可构造}。此函数是 $\operatorname{card}(D^C)$，其中 $C$ 为函数“\Term{geometric-fibre}{几何纤维}的\Term{connected-component}{连通分支}个数”；而 $C$ \Term{constructible}{可构造}（EGA IV 9.7.9），结论得证。
\end{proof}
\begin{SGAresult}{引理}{1.6}
设 $S$ 是\Term{noetherian}{诺特}\Term{scheme}{概形}，$f:X\to S$ \Term{proper}{固有}，$F$ 是 $X$ 上的\Term{constructible}{可构造}群\Term{sheaf}{层}（相应地，$A$-模\Term{sheaf}{层}），$q$ 是等于 $0$ 或 $1$ 的整数（相应地，整数 $\geq0$）。设 $S_i$（$1\leq i\leq n$）是 $S$ 的子\Term{scheme}{概形}，使得 $S=\bigcup_iS_i$，并设 $F_i$、$f_i:X_i\to S_i$ 是数据在 $S_i$ 上的“\Term{restriction}{限制}”。则若每个 $R^qf_{i*}F_i$ 都\Term{constructible}{可构造}，$R^qf_*F$ 也如此。
\end{SGAresult}
\begin{proof}
由第 XII 讲 5.1，$R^qf_*F$ 在 $S_i$ 上的\Term{restriction}{限制}\Term{isomorphism}{同构}于 $R^qf_{i*}F_i$，引理便由第 IX 讲 2.8 得到。
\end{proof}
现在进行与第 XII 讲第 8 节类似的论证。
\SourcePage{106}
\begin{SGAresult}{引理}{1.7}
设 $S$ \Term{noetherian}{诺特}，
\[
\begin{tikzcd}
X\arrow[dr,"f"']&\overline X\arrow[l,"\pi"']\arrow[d,"\overline f"]\\
&S
\end{tikzcd}
\]
是\Term{proper}{固有}\Term{morphism}{态射}的\Term{commutative-diagram}{交换图}。设 $j:U\to X$ 是一个\Term{open-subset}{开子集}，使得\Term{open-subset}{开子集} $\overline U=U\times_X\overline X$ \Term{isomorphism}{同构}地映到 $U$。设 $Y=X-U$ 配备\Term{reduced}{约化}闭子\Term{scheme}{概形}结构，$f_\circ:Y\to S$、$i:Y\to X$ 为\Term{canonical}{典范}\Term{morphism}{态射}。假定定理 1.1 对 $\overline f$、$f_\circ$ 和 $\pi$ 成立。则它也对\SGAsourcegap{原文及现有扫描补源均在“pour”后缺少\Term{object}{对象}}成立。
\end{SGAresult}
\begin{proof}
先处理 $X$ 上的 $A$-模\Term{sheaf}{层} $F$ 的情形。有\Term{exact-sequence}{正合列}
\[
0\longrightarrow j_!j^*F\longrightarrow F\longrightarrow i_*i^*F\longrightarrow0.
\]
考虑 $R^qf_*$ 的相应\Term{cohomology}{上同调}\Term{exact-sequence}{正合列}和第 IX 讲 2.6，可见只须对两端项验证 $R^qf_*$ 的\Term{constructible}{可构造}性（由第 IX 讲 2.4 (iii) 和 2.15 (i)，$i_*i^*F$ \Term{constructible}{可构造}，故 $j_!j^*F$ 也\Term{constructible}{可构造}，见第 IX 讲 2.6）。由于 $i$ 是\Term{closed-immersion}{闭浸入}，
\[
(R^qf_*)i_*(i^*F)\simeq(R^qf_{\circ*})(i^*F)
\]
（第 VIII 讲 5.6），所以由对 $f_\circ$ 的假设，得到这种情形下的\Term{constructible}{可构造}性。再记 $\overline j:\overline U\to\overline X$ 为\Term{inclusion-morphism}{包含态射}，$\overline F=\pi^*F$，则有
\[
j_!j^*F\simeq\pi_*\overline j_!\overline j^*\overline F.\tag{*}
\]
验证方法是先定义一个向右的\Term{arrow}{箭头}，再用第 XII 讲 5.1 (i) 在每个\Term{fiber}{纤维}上证明它\Term{invertible}{可逆}。对于 $Y$ 的每个\Term{geometric-point}{几何点} $\bar y$，$\overline j_!\overline j^*\overline F$ 在 $\overline X/X$ 的\Term{fiber}{纤维} $\overline X_{\bar y}$ 上的\Term{restriction}{限制}为零，因而由第 XII 讲 5.1 (iii)，对每个 $q\geq0$，$(R^q\pi_*)(\overline j_!\overline j^*\overline F)$ 在 $\bar y$ 处的\Term{fiber}{纤维}为零。另一方面，$\pi$ 在 $X-Y$ 上\Term{induction-isomorphism}{诱导同构}，故当 $q>0$ 时
\[
(R^q\pi_*)(\overline j_!\overline j^*\overline F)=0.
\]
由 \Term{leray-spectral-sequence}{Leray 谱序列}
\[
E_2^{p,q}=(R^pf_*)(R^q\pi_*)(\overline j_!\overline j^*\overline F)
\Longrightarrow R^n\overline f_*(\overline j_!\overline j^*\overline F)
\]
及 $(*)$，得到对 $q\geq0$ 的\Term{isomorphism}{同构}
\[
(R^q\overline f_*)(\overline j_!\overline j^*\overline F)\simeq(R^qf_*)(j_!j^*F).
\]
由于 $\overline j_!\overline j^*\overline F$ \Term{constructible}{可构造}（论证与 $j_!j^*F$ 相同），右端也\Term{constructible}{可构造}，从而 $(R^qf_*)(j_!j^*F)$、进而 $R^qf_*F$ \Term{constructible}{可构造}。

现在设 $F$ 是 $X$ 上的\Term{constructible}{可构造}群\Term{sheaf}{层}。以 $\pi\amalg i:(\overline X\amalg Y)\to X$ 代替 $\pi:\overline X\to X$。这样，只须证明若 $\pi:\overline X\to X$ \Term{surjectivity}{满射}，且 1.1 对 $\pi$ 和 $\overline f$ 成立，则对 $f$ 也成立。先设 $\overline F$ 是 $\overline X$ 上的\Term{constructible}{可构造}群\Term{sheaf}{层}。则 $\pi_*\overline F$ 在 $X$ 上\Term{constructible}{可构造}（1.5），
\SourcePage{107}
且有\Term{injectivity}{单射} $(R^1f_*)\pi_*\overline F\hookrightarrow(R^1\overline f_*)\overline F$。由假设右端\Term{constructible}{可构造}，故左端也\Term{constructible}{可构造}（第 IX 讲 2.9 (ii)）。由于 $\pi$ \Term{surjectivity}{满射}，有\Term{injectivity}{单射} $F\to\pi_*\pi^*F=G$。\Term{sheaf}{层} $\overline F=\pi^*F$ \Term{constructible}{可构造}，所以如上所见，$G$ 和 $R^1f_*G$ 也\Term{constructible}{可构造}。因此归结为下一个引理。
\end{proof}
\begin{SGAresult}{引理}{1.8}
设 $S$ \Term{noetherian}{诺特}，$f:X\to S$ 是\Term{proper}{固有}\Term{morphism}{态射}，$F\xhookrightarrow{u}G$ 是 $X$ 上\Term{constructible}{可构造}群\Term{sheaf}{层}的\Term{injectivity}{单射}。若 $R^1f_*G$ \Term{constructible}{可构造}，则 $R^1f_*F$ 也\Term{constructible}{可构造}。
\end{SGAresult}
\begin{proof}
设 $C=G/F$，这是一个\Term{constructible}{可构造}的 $G$-\Term{homogeneous-space}{齐性空间}\Term{sheaf}{层}（应用第 IX 讲 2.6）。考虑\Term{exact-sequence}{正合列}
\[
f_*C\xrightarrow{\delta}R^1f_*F\xrightarrow{u^1}R^1f_*G.\tag{*}
\]
为证明 $R^1f_*F$ \Term{constructible}{可构造}，由\Term{noetherian-induction}{诺特归纳}和 1.6，只须证明若 $S\ne\varnothing$，则存在 $S$ 的某个\Term{nonempty}{非空}\Term{open-subset}{开子集}，使 $R^1f_*F$ 在其上\Term{constructible}{可构造}。因此可以假定 $R^1f_*G$ \Term{locally-constant}{局部常值}；由于\Term{constructible}{可构造}性是\Term{etale}{平展}\Term{topology}{拓扑}的\Term{local-nature}{局部性质}，又可以假定 $S$ \Term{irreducible}{不可约}且 $R^1f_*G$ 常值（并\Term{constructible}{可构造}），取值 $E=\{e_1,\ldots,e_n\}$。设 $D_i$（$i=1,\ldots,n$）是\Term{section}{截面} $e_i$ 在 $u^1$ 下的\Term{inverse-image}{逆像}子\Term{sheaf}{层}。由于\Term{topos}{拓扑斯}中的“和是\Term{universal}{普遍}的”，有 $R^1f_*F=\coprod_iD_i$；故只须证明每个 $D_i$ 在某个适当的\Term{nonempty}{非空}\Term{etale}{平展} $S$-\Term{scheme}{概形} $S'$ 上\Term{constructible}{可构造}。

回顾 $R^1f_*F$ 是\Term{presheaf}{预层} $\mathscr R^1F$ 的\Term{adjunction}{伴随}\Term{sheaf}{层}，其中 $\mathscr R^1F(S')=H^1(X',F)$。因此以某个\Term{nonempty}{非空}\Term{etale}{平展} $S$-\Term{scheme}{概形} $S'$ 代替 $S$ 后，可以假定 $D_i$ 是“空”\Term{sheaf}{层}（从而\Term{constructible}{可构造}），或者存在元素 $\alpha\in H^1(X,F)$，诱导 $D_i$ 的一个\Term{section}{截面} $\bar\alpha$。在后一种情形，设
\[
f_*C^\alpha\xrightarrow{\delta}R^1f_*F^\alpha\longrightarrow R^1f_*G^\alpha
\]
是由 $(*)$ “借助 $\alpha$ \Term{twisting}{扭曲}”而得的\Term{exact-sequence}{正合列}。有
\[
R^1f_*F^\alpha\simeq R^1f_*F,\qquad R^1f_*G^\alpha\simeq R^1f_*G,
\]
且对应于 $\bar\alpha$ 的 $R^1f_*F^\alpha$ 的\Term{section}{截面}是\Term{unit-section}{单位截面}。因此，与 $D_i$ 对应的 $R^1f_*F^\alpha$ 的子\Term{sheaf}{层} $D_i^\alpha$ 是 $R^1f_*G^\alpha$ 的\Term{unit-section}{单位截面}的\Term{inverse-image}{逆像}，即 $D_i^\alpha$ 是 $f_*G^\alpha$ 在 $\delta$ 下的像。由于 $f_*G$ \Term{constructible}{可构造}（1.5），$D_i^\alpha$、从而 $D_i$ 也\Term{constructible}{可构造}（应用第 IX 讲 2.6），引理得证。
\end{proof}
\begin{SGAresult}{引理}{1.9}
设 $S$ \Term{noetherian}{诺特}，且
\[
\begin{tikzcd}
X\arrow[dr,"f"']&\overline X\arrow[l,"\pi"']\arrow[d,"\overline f"]\\
&S
\end{tikzcd}
\]
\SourcePage{108}
是\Term{proper}{固有}\Term{morphism}{态射}的\Term{commutative-diagram}{交换图}。若定理 1.1 对 $f$ 和 $\pi$ 成立，则对 $\overline f$ 也成立。
\end{SGAresult}
\begin{proof}
对于\Term{constructible}{可构造} $A$-模\Term{sheaf}{层} $\overline F$，使用 \Term{leray-spectral-sequence}{Leray 谱序列}
\[
E_2^{p,q}=(R^pf_*)(R^q\pi_*)\overline F\Longrightarrow(R^n\overline f_*)\overline F.
\]
由假设，每个 $p,q$ 的 $E_2^{p,q}$ 都\Term{constructible}{可构造}，故由第 IX 讲 2.6，\Term{abutment}{终项}也\Term{constructible}{可构造}。

设 $\overline F$ 是\Term{constructible}{可构造}群\Term{sheaf}{层}。有\Term{exact-sequence}{正合列}（第 XII 讲 3.2）
\[
0\longrightarrow(R^1f_*)\pi_*\overline F\longrightarrow(R^1\overline f_*)\overline F\longrightarrow f_*(R^1\pi_*)\overline F,\tag{*}
\]
其两端由假设和 1.5 都\Term{constructible}{可构造}。如 1.8 的证明进行：由\Term{noetherian-induction}{诺特归纳}和 1.6，只须假定 $S\ne\varnothing$，并在某个适当的\Term{nonempty}{非空}\Term{open-subset}{开子集}上证明\Term{constructible}{可构造}性。因此可假定 $f_*(R^1\pi_*)\overline F$ \Term{locally-constant}{局部常值}，甚至通过\Term{etale}{平展}\Term{localization}{局部化}，假定它常值且\Term{constructible}{可构造}（第 IX 讲 2.8），取值 $E=\{e_1,\ldots,e_n\}$。设 $D_i$ 是 $(R^1\overline f_*)\overline F$ 中 $e_i$ 的\Term{inverse-image}{逆像}子\Term{sheaf}{层}，于是 $(R^1\overline f_*)\overline F=\coprod_iD_i$。只须证明每个\Term{sheaf}{层} $D_i$ 在某个适当的\Term{nonempty}{非空}\Term{connected}{连通}\Term{etale}{平展} $S$-\Term{scheme}{概形} $S'$ 上诱导\Term{constructible}{可构造}\Term{sheaf}{层}。若 $D_i$ 不是“空\Term{sheaf}{层}”，则归结为存在元素 $\alpha\in H^1(\overline X,\overline F)$，诱导 $D_i$ 的一个\Term{section}{截面} $\bar\alpha$ 的情形。设
\[
0\longrightarrow(R^1f_*)\pi_*\overline F^\alpha
\longrightarrow(R^1\overline f_*)\overline F^\alpha
\longrightarrow f_*(R^1\pi_*)\overline F^\alpha
\]
是将 $\overline F$ 借助 $\alpha$ \Term{twisting}{扭曲}、由 $(*)$ 得到的\Term{exact-sequence}{正合列}。$(R^1\overline f_*)\overline F^\alpha$ 中对应于 $\bar\alpha$ 的\Term{section}{截面}是\Term{unit-section}{单位截面}，因此对应于 $D_i$ 的子\Term{sheaf}{层} $D_i^\alpha$ 是 $f_*(R^1\pi_*)\overline F^\alpha$ 的\Term{unit-section}{单位截面}的\Term{inverse-image}{逆像}，故
\[
D_i^\alpha\simeq(R^1f_*)\pi_*\overline F^\alpha,
\]
它确是\Term{constructible}{可构造}\Term{sheaf}{层}，证毕。
\end{proof}
\begin{SGAresult}{引理}{1.10}
只须在 $S$ 在 $\operatorname{Spec}\mathbf Z$ 上\Term{finite-type}{有限型}、$f$ \Term{projective}{射影}且\Term{relative-dimension}{相对维数} $\leq1$ 的情形验证 1.1。
\end{SGAresult}
\begin{proof}
先假定 $f$ \Term{projective}{射影}，并对 $f$ 的\Term{relative-dimension}{相对维数}归纳；以 $S$ 的一个\Term{affine}{仿射}\Term{open-subset}{开子集}代替 $S$ 后，可假定此\Term{dimension}{维数}\Term{finite}{有限}。若 $f$ 的\Term{relative-dimension}{相对维数} $\leq n$，$n\geq2$，则在 $S$ 上\Term{local-scheme}{局部}可找到\Term{finite-morphism}{有限态射} $\pi:X\to\mathbf P_S^n$，由 1.9 和第 IX 讲 2.15 (i)，归结为 $f$ 是 $\mathbf P_S^n$ 的\Term{structural-morphism}{结构态射}的情形。考虑\Term{commutative-diagram}{交换图}
\SourcePage{109}
\[
\begin{tikzcd}
&\overline X\arrow[dl,"\pi"']\arrow[dr,"a"]\arrow[dd,"\overline f" description]&\\
\mathbf P_S^n\arrow[dr,"f"']&&\mathbf P_S^1\arrow[dl,"b"]\\
&S&
\end{tikzcd}
\]
此图在第 XII 讲 8.3 的证明中已经出现，其中 $\overline X$ 由 $\mathbf P_S^n$ 沿闭子\Term{scheme}{概形} $Y\simeq\mathbf P_S^{n-2}$ 吹起而得。$b$ 的\Term{relative-dimension}{相对维数} $\leq1$，$a$ 的\Term{relative-dimension}{相对维数} $\leq n-1$，故由归纳假设，定理对 $a$、$b$ 成立，由 1.9 对 $\overline f$ 也成立。此外，$\pi$ 的\Term{dimension}{维数} $\leq1$，$f_\circ:Y\to S$ 的\Term{dimension}{维数} $\leq n-1$。因此由 1.7，1.1 对 $f$ 成立，从而对每个\Term{projective}{射影}\Term{morphism}{态射}成立。

现在设 $f$ 是任意\Term{proper}{固有}\Term{morphism}{态射}，对 $X$ 作\Term{noetherian-induction}{诺特归纳}：可以假定定理对 $X$ 的每个真闭子\Term{scheme}{概形}成立，并假定 $X\ne\varnothing$。取 Chow 引理（第 XII 讲 7.1）所给类型的图
\[
\begin{tikzcd}
X\arrow[dr,"f"']&\overline X\arrow[l,"\pi"']\arrow[d,"\overline f"]\\
&S.
\end{tikzcd}
\]
把 1.7 用于此图，得到 $f$ 的 1.1，引理得证。
\end{proof}
\textbf{1.11。}现在可以完成 1.1 的证明。若 $F$ 是\Term{constructible}{可构造}群\Term{sheaf}{层}，可以找到\Term{injectivity}{单射} $F\hookrightarrow G=\prod_i\pi_{i*}C_i$，其中 $\pi_i$ \Term{finite}{有限}，$C_i$ 常值且\Term{constructible}{可构造}（第 IX 讲 2.15）；由 1.8，可以假定 $F$ 常值。同样，若 $F$ 是\Term{constructible}{可构造} $A$-模\Term{sheaf}{层}，则应用第 IX 讲 2.15，可找到 $F$ 的一个\Term{resolution}{分解}
\[
F\longrightarrow G^\bullet,\qquad G^\bullet=\{0\longrightarrow G^0\longrightarrow G^1\longrightarrow G^2\longrightarrow\cdots\},
\]
使每个 $G^\nu$ 都具有 $G^\nu=\prod_i\eta_{i*}C_i$ 的\Term{form}{形式}。由\Term{resolution}{分解} $G^\bullet$ 的\Term{spectral-sequence}{谱序列}
\[
E_2^{p,q}=H^q((R^qf_*)G^\bullet)\Longrightarrow(R^nf_*)F,
\]
归结为 $F$ 常值且\Term{constructible}{可构造}、取值为 $M$ 的情形。

回顾 $A$ 是 $\mathbf Z/n$-\Term{algebra}{代数}。设 $\ell$ 是整除 $n$ 的素数。有\Term{exact-sequence}{正合列}
\[
0\longrightarrow\ell M\longrightarrow M\longrightarrow M/\ell M\longrightarrow0,
\]
因此又只须对 $\ell M$ 和 $M/\ell M$ 证明 1.1。
\SourcePage{110}
由简单归纳，便归结为 $A$ 是 $\mathbf Z/\ell$-\Term{algebra}{代数}、$\ell\in\mathbf P$、$M$ 是\Term{finite-type}{有限型} $A$-模的情形。于是 $M$ 是 $\mathbf Z/\ell$ 上的\Term{vector-space}{向量空间}，从而\Term{canonical}{典范}\Term{morphism}{态射}
\[
H^q(Y,\mathbf Z/\ell)\otimes_{\mathbf Z/\ell}M\longrightarrow H^q(Y,M)
\]
对任意\Term{quasi-compact}{拟紧}且\Term{quasi-separated}{拟分离}的 $Y$ 都\Term{bijection}{双射}（事实上，当 $M$ \Term{finite-rank}{有限秩}时这是平凡的，而两端都与\Term{inductive-limit}{归纳极限}可交换，见第 VI 讲 1.6）。故
\[
[R^qf_*(\mathbf Z/\ell)]\otimes_{\mathbf Z/\ell}M\xrightarrow{\sim}R^qf_*(M).
\]
所以显然只须证明 $R^qf_*(\mathbf Z/\ell)$ 作为 $\mathbf Z/\ell$-模\Term{constructible}{可构造}，即归结为 $A=\mathbf Z/\ell$ 的情形，下文即作此假定。

注意，$f$ 的\Term{relative-dimension}{相对维数} $\leq0$（即 $f$ \Term{finite}{有限}）、$F$ 任意的情形，现在由第 VIII 讲 5.6 与 1.5 得到。

对 $S$ 作\Term{noetherian-induction}{诺特归纳}。由 1.6，可以假定 $S\ne\varnothing$，并允许以任意\Term{nonempty}{非空}\Term{open-subset}{开子集}代替 $S$。还可以假定 $S$ 整，\Term{generic-point}{泛点}为 $s$。\Term{fiber}{纤维} $X_s$ 是\Term{dimension}{维数} $\leq1$ 的\Term{algebra}{代数}\Term{scheme}{概形}，故存在 $k(s)=K$ 的一个\Term{radicial}{根性}\Term{group-extension}{扩张} $K'$，使 $(X_s\otimes_KK')_{\mathrm{red}}$ 的\Term{normalization}{正规化} $Y$ 在 $\operatorname{Spec}K'$ 上\Term{smooth}{光滑}\footnote{见 EGA IV 17.15.14。}。以一个\Term{nonempty}{非空}\Term{open-subset}{开子集}代替 $S$ 后，可以假定存在\Term{finite}{有限}\Term{radicial}{根性}\Term{epimorphism}{满态射} $S'\to S$，使得 $S'_s$ 与 $\operatorname{Spec}K'$ 在 $K$ 上\Term{isomorphism}{同构}，以及一个\Term{finite}{有限}\Term{epimorphism}{满态射}\SGAsourcegap{此处两版均仅有“un morphisme fini surjectif”这一残句，未给出该\Term{morphism}{态射}的源、目标及名称}、一个\Term{smooth}{光滑}且\Term{projective}{射影}的\Term{morphism}{态射} $\overline X'\to S'$ 和一个\Term{finite}{有限}满的 $S'$-\Term{morphism}{态射}
\[
\Pi:\overline X'\longrightarrow X'=X\times_SS',
\]
其在 $S$ 的\Term{generic-point}{泛点} $s$ 处的\Term{fiber}{纤维}上诱导 $Y\to X_s\times_SK'$（EGA IV 9.6.1）。由第 VIII 讲 1.1，作\Term{radicial}{根性}\Term{group-extension}{扩张}没有影响。因此可以用 $S'$ 代替 $S$、用 $X'$ 代替 $X$，即假定 $f$ 包含在图
\[
\begin{tikzcd}
X\arrow[dr,"f"']&\overline X\arrow[l,"\pi"']\arrow[d,"\overline f"]\\
&S
\end{tikzcd}
\]
中，其中 $\pi$ \Term{finite}{有限}且\Term{surjectivity}{满射}，$\overline f$ \Term{smooth}{光滑}。而且由构造，可以假定存在 $X$ 的一个闭子\Term{scheme}{概形} $Y$，其\Term{generic-fiber}{泛纤维}\Term{finite}{有限}，且 $\pi$ 在 $X-Y$ 上\Term{induction-isomorphism}{诱导同构}。再次以\Term{nonempty}{非空}\Term{open-subset}{开子集}代替 $S$，可以进一步归结为 $f_\circ:Y\to S$ 是\Term{finite-morphism}{有限态射}的情形。

由于 $F$ 常值，$\pi^*F$ 也常值；由于 $\pi$ \Term{surjectivity}{满射}，$F\to\pi_*\pi^*F$ \Term{injectivity}{单射}。若 $F$ 是群\Term{sheaf}{层}，只须证明
\SourcePage{111}
\[
(R^1f_*)(\pi_*\pi^*F)\simeq(R^1\overline f_*)(\pi^*F)
\]
\Term{constructible}{可构造}（1.8）。若 $F$ 是 $\mathbf Z/\ell$-模\Term{sheaf}{层}，有\Term{exact-sequence}{正合列}
\[
0\longrightarrow F\longrightarrow\pi_*\pi^*F\longrightarrow C\longrightarrow0,
\]
其中 $C$ 集中在 $Y$ 上，因为 $\pi$ 在 $X-Y$ 上\Term{induction-isomorphism}{诱导同构}。因此 $R^qf_*C=R^qf_{\circ*}C$ \Term{constructible}{可构造}（$q>0$ 时为零），从而仍只须证明
\[
(R^qf_*)(\pi_*\pi^*F)=(R^q\overline f_*)(\pi^*F)
\]
\Term{constructible}{可构造}。以 $\overline f$ 代替 $f$、以 $\pi^*F$ 代替 $F$，便归结为 $F$ 常值且\Term{constructible}{可构造}、$f$ \Term{smooth}{光滑}的情形。

先处理\Term{constant-group}{常值群}\Term{sheaf}{层} $F=G_X$。$(R^1f_*)G_X$ 的所有\Term{fiber}{纤维}都\Term{finite}{有限}；事实上，应用第 XII 讲 5.2 (ii)，为证明此事，可归结为 $S$ 是\Term{separably-closed-field}{可分闭域}的谱、$X$ 在 $S$ 上\Term{smooth}{光滑}且\Term{dimension}{维数} $\leq1$ 的情形。此时结果是已知的（SGA I 第 X 讲 2.6）。因此只须证明\Term{specialization-morphism}{特化态射}\Term{injectivity}{单射}（第 IX 讲 2.13 (ii)）。为此可归结为 $S$ \Term{strictly-local}{严格局部}的情形；由\Term{base-change-theorem}{基变换定理}第 XII 讲 5.1 (ii)，甚至可以用一个\Term{strictly-local}{严格局部}的\Term{discrete-valuation-ring}{离散赋值环}的谱代替 $S$。又可以假定 $X$ \Term{connected}{连通}。由于 $H^1(X,G)$ 分类 $X$ 的群为 $G$ 的\Term{principal-galois-covering}{主 Galois 覆盖}，只须验证若 $X'\to X$ 是\Term{connected}{连通}\Term{etale-covering}{平展覆盖}，则\Term{geometric-fibre}{几何纤维}也\Term{connected}{连通}。这容易由 $f$ \Term{smooth}{光滑}且\Term{proper}{固有}、\Term{geometric-fibre}{几何纤维}\Term{connected}{连通}得到。

现在处理 $F=(\mathbf Z/\ell)_X$ 的情形。$q$ 等于 $0$ 或 $1$ 的结果已知（1.5 以及上文）。由第 X 讲 4.3 和第 XII 讲 5.3 bis，只剩 $q=2$ 的情形；而且，若 $S$ 的\Term{characteristic}{特征}是 $\ell$，则 $R^2f_*(\mathbf Z/\ell)=0$。因此，以适当的\Term{nonempty}{非空}\Term{open-subset}{开子集}代替 $S$ 后，归结为 $\ell$ 在 $S$ 上\Term{invertible}{可逆}的情形。此时由第 XII 讲 5.2 (ii) 和第 X 讲 4.7，$R^2f_*(\mathbf Z/\ell)$ 在位于 $s\in S$ 之上的\Term{geometric-point}{几何点} $\bar s$ 处的\Term{fiber}{纤维}\Term{isomorphism}{同构}于 $(\mathbf Z/\ell)^{c(s)}$，其中 $c(s)$ 是\Term{geometric-fibre}{几何纤维} $X_{\bar s}$ 的一维\Term{irreducible-component}{不可约分支}个数。而 $c$ 是 $S$ 上的\Term{constructible}{可构造}函数（EGA IV 9.7.9），结论得证（第 IX 讲 2.13 (iii)）。

\textbf{注记 1.11。}利用\Term{proper-support}{固有支撑}\Term{cohomology}{上同调}的\Term{form}{形式}体系，1.1 的证明会大为简化；这个\Term{form}{形式}体系将在第 XVII 讲中作为\Term{proper}{固有}\Term{morphism}{态射}的\Term{base-change-theorem}{基变换定理}第 XII 讲 5.1 的直接推论建立。另一方面，1.1 可推广到如下情形：在 $Y$ 上给定任意挠环\Term{sheaf}{层} $A$，并在 $X$ 上取一个 $f^*(A)$-模\Term{complex}{复形} $K^\bullet$，满足某种“\Term{constructible}{可构造}性”条件；在 1.1 的特殊情形中，这个条件可简单地表为：$K^\bullet$ 的\Term{cohomology}{上同调}\Term{sheaf}{层} $\mathcal H^i(K^\bullet)$ 是\Term{constructible}{可构造} $A$-模，并在 $i$ 充分大时为零。

\SGAsection{2}{维数的一种变式}
\textbf{2.1。}设 $X$ 是\Term{scheme}{概形}。对于 $x\in X$，记
\[
d(x)=\dim\overline{\{x\}},\tag{2.1.1}
\]
\SourcePage{112}
其中 $\overline{\{x\}}$ 是 $\{x\}$ 的\Term{closure}{闭包}。若 $F$ 是 $X$ 上的阿贝尔\Term{sheaf}{层}，置
\[
d(F)=\sup\{d(x)\mid x\in X,\ F_{\bar x}\ne0\}.\tag{2.1.2}
\]
在\Term{local-scheme}{局部}情形需要一个修正的概念。
\begin{SGAresult}{定义}{2.2}
设 $Y$ 是一个\Term{universal}{普遍}悬链的\Term{noetherian-local-ring}{诺特局部环}的谱（EGA IV 5.6.2），$f:X\to Y$ 是\Term{finite-type}{有限型}\Term{morphism}{态射}。设 $x\in X$，$y=f(x)$。置
\[
\delta(x)=\dim\overline{\{y\}}+\operatorname{trdeg}_{k(y)}k(x).
\]
若 $F$ 是 $X$ 上的阿贝尔\Term{sheaf}{层}，置
\[
\delta(F)=\sup\{\delta(x)\mid x\in X,\ F_{\bar x}\ne0\}.
\]
\end{SGAresult}
\begin{SGAresult}{命题}{2.3}
(i) 设 $i:X'\to X$ 是一个\Term{immersion}{浸入}，$x'\in X'$。则 $\delta(i(x'))=\delta(x')$。

(ii) 设 $X_\circ$ 是 $X/Y$ 的\Term{closed-fiber}{闭纤维}。有
\[
\delta(x)=\begin{cases}
d(x)&\text{若 }\overline{\{x\}}\cap X_\circ\ne\varnothing,\\
d(x)+1&\text{若 }\overline{\{x\}}\cap X_\circ=\varnothing.
\end{cases}
\]
(iii) 若 $f$ \Term{proper}{固有}，则 $\delta(x)=d(x)$。
\end{SGAresult}
\begin{proof}
(i) 由定义显然成立。当 $f$ \Term{proper}{固有}时，$\overline{\{x\}}\cap X_\circ$ \Term{nonempty}{非空}，因此 (iii) 是 (ii) 的推论。对于 (ii)，由 (i) 可将 $X$、$Y$ 分别替换为 $\overline{\{x\}}$、$\overline{\{y\}}$（配备诱导的\Term{reduced-structure}{约化结构}）。设 $x'\in X$，$y'=f(x')$。有公式（EGA IV 5.6.1.1）
\[
\dim\mathcal O_{X,x'}=\dim\mathcal O_{Y,y'}+\operatorname{trdeg}_{k(y)}k(x)-\operatorname{trdeg}_{k(y')}k(x').
\]
假定 $X_\circ\ne\varnothing$。显然，$\sup_{x'}\{\dim\mathcal O_{X,x'}\}=\dim X$ 在取 $x'$ 为 $X_\circ$ 的一个\Term{closed-point}{闭点}时取得，故
\[
\dim X=\dim\mathcal O_{X,x'}=\dim Y+\operatorname{trdeg}_{k(y)}k(x)=\delta(x).
\]
若 $X_\circ=\varnothing$，则显然在取 $x'$ 为其\Term{fiber}{纤维} $f^{-1}(y')$ 中的\Term{closed-point}{闭点}、且满足 $\dim\overline{\{y'\}}=1$ 时取得\Term{supremum}{上确界}。这种点存在，因为满足 $\dim\overline{\{y'\}}=1$ 的 $y'\in Y$ 构成一个很\Term{dense}{稠密}的\Term{set}{集合}（EGA IV 10.1.3、10.5.9、1.8.4）。因此有
\[
\begin{aligned}
\dim X&=\dim\mathcal O_{X,x'}=\dim\mathcal O_{Y,y'}+\operatorname{trdeg}_{k(y)}k(x)\\
&=\dim Y-1+\operatorname{trdeg}_{k(y)}k(x)=\delta(x)=1,
\end{aligned}
\]
结论得证。
\end{proof}
\SourcePage{113}
\begin{SGAresult}{命题}{2.4}
设 $f:X\to Y$ 是域 $k$ 的谱上\Term{finite-type}{有限型}\Term{scheme}{概形}的\Term{morphism}{态射}。设 $x\in X$，$y=f(x)$，$y_\circ$ 是 $y$ 的一个\Term{specialization}{特化}。设 $f':X'\to Y'$ 是 $f$ 在位于 $y_\circ$ 之上的\Term{geometric-point}{几何点} $\bar y_\circ$ 处的\Term{strict-localization}{严格局部化}，$x'\in X'$ 是 $x$ 上方的一点。则有
\[
\delta(x')=d(x)-d(y_\circ).
\]
\end{SGAresult}
\begin{proof}
有
\[
\begin{aligned}
d(x)&=d(y)+\operatorname{trdeg}_{k(y)}k(x),\\
\delta(x')&=d(y')+\operatorname{trdeg}_{k(y')}k(x'),
\end{aligned}
\]
其中 $y'=f'(x')$。此外，
\[
\operatorname{trdeg}_{k(y)}k(x)=\operatorname{trdeg}_{k(y')}k(x'),
\]
因为\Term{group-extension}{扩张} $k(x')/k(x)$ 和 $k(y')/k(y)$ 都是\Term{algebra}{代数}\Term{group-extension}{扩张}。因此只须证明
\[
d(y')=d(y)-d(y_\circ).
\]
可以假定 $y$ 在 $Y$ 中极大，故 $y'$ 在 $Y'$ 中极大；于是
\[
\begin{aligned}
d(y')&=\dim\mathcal O_{Y',\bar y_\circ}\\
&=\dim\mathcal O_{Y,y_\circ}&&\text{（EGA IV 6.1.3）}\\
&=\dim Y-d(y_\circ)&&\text{（EGA IV 5.2.3.1）}\\
&=d(y)-d(y_\circ),
\end{aligned}
\]
结论得证。
\end{proof}
\SGAsection{3}{仿射代数概形的上同调维数}
\begin{SGAresult}{定理}{3.1 ($d$)}
设 $f:X\to Y$ 是域 $k$ 上\Term{finite-type}{有限型}\Term{scheme}{概形}的\Term{affine-morphism}{仿射态射}。设 $F$ 是 $X$ 上的\Term{torsion-sheaf}{挠层}，满足 $d(F)\leq d$（见 2.1）。则 $d(R^qf_*F)\leq d-q$。
\end{SGAresult}
特别地，取 $Y=\operatorname{Spec}k$：
\begin{SGAresult}{推论}{3.2}
设 $X$ 是\Term{separably-closed-field}{可分闭域} $k$ 上的\Term{finite-type}{有限型}\Term{affine}{仿射}\Term{scheme}{概形}。则
\[
\operatorname{cd}X\leq\dim X.
\]
\end{SGAresult}
\begin{SGAresult}{推论}{3.3}
设 $X\subset\mathbf P_k^N$ 是\Term{separably-closed-field}{可分闭域} $k$ 上的 $n$ 维\Term{projective}{射影}\Term{scheme}{概形}，$Y=H\cap X$ 是 $X$ 的一个“\Term{hyperplane-section}{超平面截面}”。则对 $X$ 上的每个\Term{torsion-sheaf}{挠层} $F$，\Term{canonical-homomorphism}{典范同态}
\[
H_Y^q(X,F)\longrightarrow H^q(X,F)
\]
在 $q>n+1$ 时\Term{bijection}{双射}，在 $q=n+1$ 时\Term{surjectivity}{满射}。
\end{SGAresult}
\SourcePage{114}
考虑到 $X-Y=U$ \Term{affine}{仿射}、\Term{dimension}{维数} $\leq n$，此推论由 2.2 和\Term{closed-subset}{闭子集}的\Term{exact-sequence}{正合列}（第 IV 讲 3.7.1）立即得到。注意 3.3 推广了 Lefschetz 的一个\Term{hyperplane-section}{超平面截面}定理：事实上，若 $X$ 在 $k$ 上\Term{smooth}{光滑}、$F$ \Term{locally-constant}{局部常值}且其阶与\Term{characteristic}{特征}互素，则利用\Term{relative-cohomology}{相对上同调}\Term{purity-theorem}{纯性定理}（第 XVI 讲第 3 节）及\Term{spectral-sequence}{谱序列}（\SGAsourcegap{原文与现有扫描补源的\Term{spectral-sequence}{谱序列}引用编号均为空白}），可以显式计算\Term{cohomology}{上同调} $H_Y^q(X,F)$，得到对于所有 $q$，
\[
H_Y^q(X,\mathbf Z/n)\simeq H^{q-2}(Y,\mu_n^{-1});
\]
由 3.3 得到的\Term{homomorphism}{同态} $H^{q-2}(Y,\mu_n^{-1})\to H^q(X,\mathbf Z/n)$，正是以后将要研究的“\Term{gysin-homomorphism}{Gysin 同态}”\footnote{关于 Lefschetz 型定理的系统研究，见 SGA 2 第 XIV 讲。}。

定理 3.1 ($d$) \Term{equivalence}{等价}于下面的结果。
\begin{SGAresult}{推论}{3.4 ($d$)}
设 $Y'$ 是\Term{strictly-local}{严格局部}\Term{scheme}{概形}，它是域上某个\Term{algebra}{代数}\Term{scheme}{概形}的\Term{strict-localization}{严格局部化}。设 $f':X'\to Y'$ 是\Term{finite-type}{有限型}\Term{affine-morphism}{仿射态射}，$F'$ 是 $X'$ 上的\Term{torsion-sheaf}{挠层}，满足 $\delta(F')\leq d$（见 2.2）。则当 $q>d$ 时，$H^q(X',F')=0$。
\end{SGAresult}
事实上，假定 3.4 ($d$) 成立。设 $f:X\to Y$ 是 $k$ 上\Term{algebra}{代数}\Term{scheme}{概形}的\Term{affine-morphism}{仿射态射}，$F$ 是 $X$ 上的\Term{torsion-sheaf}{挠层}，满足 $d(F)\leq d$。设 $y$ 是 $Y$ 的一点，$\xi$ 是 $y$ 上方的一个\Term{geometric-point}{几何点}。为证明 2.1 ($d$)，须证明若 $d(y)>d-q$，则\Term{fiber}{纤维} $(R^qf_*F)_\xi$ 为零。

设 $f':X'\to Y'$ 是 $f$ 在 $\xi$ 处的\Term{strict-localization}{严格局部化}，即 $Y'$ 是 $Y$ 在 $\xi$ 处的\Term{strict-localization}{严格局部化}，$X'=X\times_YY'$。设 $F'$ 是 $F$ 在 $X$ 上诱导的\Term{sheaf}{层}，从而 $(R^qf_*F)_\xi=H^q(X',F')$（第 VIII 讲 5.2）。有
\[
\delta(F')\leq d(F)-d(y)\leq d-d(y)\qquad\text{（2.4）},
\]
应用 3.4 ($d$)，得 $q>d-d(y)$ 时 $H^q(X',F')=0$，结论得证。

反过来，假定 2.1 ($d$) 成立，并设 $F'$、$f':X'\to Y'$ 如 3.4 ($d$) 的陈述。由于\Term{cohomology}{上同调}与\Term{inductive-limit}{归纳极限}可交换（第 VII 讲 3.3），且 $F'$ 是其\Term{constructible}{可构造}子\Term{sheaf}{层}的\Term{inductive-limit}{归纳极限}（第 IX 讲 2.9），可以假定 $F'$ \Term{constructible}{可构造}。$Y'$ 是某个\Term{algebra}{代数}\Term{scheme}{概形} $Y$ 在位于\Term{closed-point}{闭点} $y_\circ$ 之上的\Term{geometric-point}{几何点} $\xi$ 处的\Term{strict-localization}{严格局部化}（第 X 讲 3.3）。写成 $Y'=\varprojlim_\alpha Y_\alpha$，其中 $(Y_\alpha)$ 是 $Y$ 上的\Term{affine}{仿射}\Term{etale}{平展}\Term{scheme}{概形}的\Term{filtered}{滤过}\Term{projective-system}{逆系}（第 VIII 讲 4.5）。由第 IX 讲 2.7.4，$\alpha$ 充分大时，$F'$ 由 $X=X\times_YY_K$ 上的一个\Term{constructible}{可构造}\Term{torsion-sheaf}{挠层}诱导；利用 2.4 与 EGA IV 9.5.5，立即可见还可以假定 $d(F_\alpha)=\delta(F')\leq d$。以 $Y_\alpha$ 代替 $Y$，便归结为 $F'$ 由 $X$ 上的\Term{sheaf}{层} $F$ 诱导、且 $d(F)\leq d$ 的情形。由 2.1 ($d$)，当 $q>D$ 时有 $R^qf_*F=0$，故特别地
\[
(R^qf_*F)_\xi=H^q(X',F')=0\qquad\text{若 }q>d,
\]
\SourcePage{115}
从而得到 3.4 ($d$)。作为 3.4 ($d$) 的特殊情形，有：
\begin{SGAresult}{推论}{3.5}
设 $Y'$ 是\Term{dimension}{维数} $\leq d$ 的\Term{strictly-local}{严格局部}\Term{scheme}{概形}，它是域 $k$ 上某个\Term{algebra}{代数}\Term{scheme}{概形}的\Term{strict-localization}{严格局部化}。设 $U\subset Y'$ 是\Term{affine}{仿射}\Term{open-subset}{开子集}。则 $\operatorname{cd}U\leq d$。
\end{SGAresult}
\textbf{注记 3.6。}看来很有可能，3.4、从而 3.5 对任意\Term{noetherian}{诺特}\Term{local-scheme}{局部}\Term{scheme}{概形} $Y'$ 都成立（未必是\Term{algebra}{代数}\Term{scheme}{概形}的\Term{strict-localization}{严格局部化}），至少在 $Y'$ \Term{excellent}{优良}（EGA IV 7.8）时如此。在进一步假定 $Y$ 的\Term{characteristic}{特征}为零时，第 XX 讲第 6 节将证明这一点。
\SGAsection{4}{定理 3.1 的证明}
证明对 $d$ 归纳。
\begin{SGAresult}{引理}{4.1}
3.1 (0) 和 3.1 (1) 成立。
\end{SGAresult}
\begin{proof}
$d=0$ 时平凡，此时 $\operatorname{Supp}F$ 在 $k$ 上\Term{finite}{有限}，应用第 VIII 讲 5.5 即可。设 $F$、$f:X\to Y$ 如 3.1 所述，并假定 $d(F)\leq1$。以 $F$ 的某个\Term{constructible}{可构造}子\Term{sheaf}{层}代替 $F$（第 IX 讲 2.9 与第 VII 讲 3.3），便归结为 $F$ \Term{constructible}{可构造}的情形。此时 $F$ 的\Term{support}{支撑}包含在 $X$ 的一个\Term{dimension}{维数} $\leq1$ 的闭子\Term{scheme}{概形}中（第 IX 讲 2.3 (ii)），并且可以用此闭子\Term{scheme}{概形}代替 $X$，即可以假定 $\dim X\leq1$。由第 IX 讲 5.7，\Term{separably-closed-field}{可分闭域}上的一维\Term{affine-algebra}{仿射代数}\Term{scheme}{概形}的\Term{cohomological-dimension}{上同调维数} $\leq1$。而且只须处理 $k$ \Term{separably-closed}{可分闭}的情形：事实上，设 $\overline f:\overline X\to\overline Y$ 是由 $f$ 经\Term{base-change}{基变换} $\operatorname{Spec}\bar k\to\operatorname{Spec}k$ 得到的\Term{morphism}{态射}（$\bar k$ 是 $k$ 的\Term{separable-closure}{可分闭包}）。对 $Y$ 的一个\Term{geometric-point}{几何点} $\xi$，$f$ 在 $\xi$ 处的\Term{strict-localization}{严格局部化} $f':X'\to Y'$ 可等同于 $\overline f$ 的一个\Term{strict-localization}{严格局部化}，因此\Term{fiber}{纤维} $(R^qf_*F)_\xi$ 可等同于 $R^q\overline f_*\overline F$ 的一个\Term{fiber}{纤维}。由于 $\overline f$ \Term{affine}{仿射}，得到 $q>1$ 时 $R^qf_*F=0$。而且显然 $d(f_*F)\leq1$。对于 $q=1$，注意 $\dim X\leq1$ 的假设意味着存在 $Y$ 的\Term{finite}{有限}\Term{closed-subset}{闭子集} $Z$，使 $f$ 在 $Y-Z$ 上\Term{finite}{有限}。因此 $R^1f_*F$ 在 $Y-Z$ 上为零（第 VIII 讲 5.6），故 $d(R^1f_*F)\leq1$，引理的证明完成。
\end{proof}
设 $d\geq2$，现在假定对 $d'<d$，3.1 ($d'$)（\Term{equivalence}{等价}地，3.4 ($d'$)）成立。
\begin{SGAresult}{引理}{4.2}
为在所有情形证明 3.4 ($d$)，只须处理 $X'=\mathbf E_{Y'}^1$ 是 $Y'$ 上的一维\Term{affine-space}{仿射空间}、$f$ 是\Term{structural-morphism}{结构态射}的情形。
\end{SGAresult}
\begin{proof}
显然可假定对适当的 $N$ 有 $X'=\mathbf E_{Y'}^N$，因为可以把 $X'$ \Term{embedding}{嵌入}某个 $\mathbf E_{Y'}^N$。设 $N>1$，并假定对所有 $Y'$ 和 $r<N$ 的 $\mathbf E_{Y'}^r$ 已知此结果。设 $F'$ 是 $\mathbf E_{Y'}^N$ 上的\Term{torsion-sheaf}{挠层}，满足 $d(F')\leq d$。
\SourcePage{116}
考虑图
\[
\begin{tikzcd}
\mathbf E_{Y'}^N\arrow[rr,"g'"]\arrow[dr,"f'"']&&\mathbf E_{Y'}^{N-1}\arrow[dl,"h'"]\\
&Y'&
\end{tikzcd}
\]
其中 $g'$ 和 $h'$ 为\Term{canonical-projection}{典范投影}，因此 $\mathbf E_{Y'}^N$ \Term{isomorphism}{同构}于 $\mathbf E_{Y'}^{N-1}$ 上的一维\Term{affine-space}{仿射空间}。显然，$\mathbf E_{Y'}^{N-1}$ 的每个\Term{strict-localization}{严格局部化}环也都是\Term{algebra}{代数}\Term{scheme}{概形}的\Term{strict-localization}{严格局部化}，故可将归纳假设用于 $R^qg'_*F'$ 的\Term{fiber}{纤维}，得到
\[
\delta(R^qg'_*F')\leq d-q.\tag{*}
\]
此外，由对 $N$ 的归纳假设，3.4 ($d'$) 对 $h'$ 成立，其中 $d'\leq d$。考虑 \Term{leray-spectral-sequence}{Leray 谱序列}
\[
E_2^{p,q}=H^p(\mathbf E_{Y'}^{N-1},R^qg'_*(F'))\Longrightarrow H^n(\mathbf E_{Y'}^N,F').
\]
将 $(*)$ 与 3.4 ($d'$) 用于\Term{morphism}{态射} $h'$，得到当 $p>d-q$ 时 $E_2^{p,q}=0$。因此当 $n>d$ 时\Term{abutment}{终项}为零，引理得证。
\end{proof}
考虑下面的附加陈述，它是 3.5 的特殊情形。
\begin{SGAresult}{陈述}{4.3 ($d$)}
设 $X'$ \Term{strictly-local}{严格局部}，\Term{dimension}{维数} $\leq d$，并且是域上某个\Term{algebra}{代数}\Term{scheme}{概形}的\Term{strict-localization}{严格局部化}。设 $f\in\Gamma(X',\mathcal O_{X'})$，$U\subset X'$ 为\Term{open-subset}{开子集} $U=X'-V(f)$。则 $\operatorname{cd}U\leq d$。
\end{SGAresult}
\begin{SGAresult}{引理}{4.4}
4.3 ($d$) 蕴含 3.4 ($d$)。
\end{SGAresult}
\begin{proof}
考虑图
\[
\begin{tikzcd}
\mathbf E_{Y'}^1\arrow[rr,"i"]\arrow[dr,"f'"']&&\mathbf P_{Y'}^1\arrow[dl,"g'"]\\
&Y'&
\end{tikzcd}
\]
其中 $Y'$ 如 2.6 的陈述，$\mathbf P_{Y'}^1$ 是由 $\mathbf E^1$ “添加无穷远\Term{section}{截面} $Y'_\infty$”而得的\Term{projective-space}{射影空间}。设 $F'$ 是 $\mathbf E^1$ 上的\Term{torsion-sheaf}{挠层}，满足 $d(\operatorname{Supp}F')\leq d$。有\Term{spectral-sequence}{谱序列}
\[
H^p(\mathbf P^1,R^qi_*F')\Longrightarrow H^n(\mathbf E^1,F').
\]
当 $q>0$ 时，$R^qi_*F'$ 集中在 $Y'_\infty$ 上，而 $Y'_\infty$ 在 $Y'$ 上\Term{isomorphism}{同构}于 $Y'$。由于 $Y'$ \Term{strictly-local}{严格局部}，当 $p,q>0$ 时有 $H^p(\mathbf P^1,R^qi_*F')=0$。此外，由\Term{proper}{固有}\Term{morphism}{态射}的\Term{base-change-theorem}{基变换定理}（第 XII 讲 5.1），
\SourcePage{117}
$\mathbf P^1$ 上\Term{torsion-sheaf}{挠层}的\Term{cohomology}{上同调}可以在\Term{closed-fiber}{闭纤维}上计算，而\Term{closed-fiber}{闭纤维}是\Term{cohomological-dimension}{上同调维数} $\leq2$ 的\Term{scheme}{概形}（第 X 讲 4.3）。因此 $p>2$ 时 $H^p(\mathbf P^1,i_*F')=0$。要证明 $n>d$ 时 $H^n(\mathbf E^1,F')=0$；由于 $d\geq2$，只须证明 $q>d$ 时 $H^0(\mathbf P^1,R^qi_*F')=0$。这个群也\Term{isomorphism}{同构}于 $H^0(Y'_\infty,R^1i_*(F'))$；由第 VIII 讲 4.6，它\Term{isomorphism}{同构}于 $R^qi_*F'$ 在点 $Q$ 处的\Term{fiber}{纤维}，其中 $Q$ 是\Term{closed-fiber}{闭纤维}中 $\mathbf P^1$ 的无穷远点。

只须处理 $F'$ \Term{constructible}{可构造}的情形（第 VII 讲 3.3 与第 IX 讲 2.9）。于是由 $\delta(\operatorname{Supp}F')\leq d$，存在 $\mathbf P^1$ 的一个闭子\Term{scheme}{概形} $X'$，满足 $\dim X'\leq d$，且包含 $F'$ 的\Term{support}{支撑}。设 $\widetilde X'$ 是 $X'$ 在\Term{geometric-point}{几何点} $Q$ 处的\Term{localization}{局部化}，置
\[
\widetilde U=\widetilde X'-\widetilde X'\times_{\mathbf P^1}Y'_\infty,
\]
并设 $\widetilde F$ 是 $F'$ 经\Term{morphism}{态射} $\widetilde U\to\mathbf P^1$ 在 $\widetilde U$ 上诱导的\Term{sheaf}{层}。那么 $R^qi_*F$ 在 $Q$ 处的\Term{fiber}{纤维}正是 $H^p(\widetilde U,\widetilde F)$（第 VIII 讲 5.2）；把 4.3 ($d$) 用于\Term{strictly-local}{严格局部}\Term{scheme}{概形} $X'$ 及\Term{open-subset}{开集} $\widetilde U$ 即可。
\end{proof}
下面的引理与 4.4 一起将完成证明。
\begin{SGAresult}{引理}{4.5}
对 $d'<d$ 的 3.4 ($d'$) 蕴含对 $d'\leq d$ 的 4.3 ($d'$)。
\end{SGAresult}
\begin{proof}
假定对 $d'<d$，3.4 ($d'$) 成立。通过归纳，可以假定对 $d'<d$，4.3 ($d'$) 成立。设数据 $X'$、$f$、$U$ 如 4.3 ($d$) 的陈述。则 $X$ 是域 $k$ 上某个\Term{algebra}{代数}\Term{scheme}{概形} $X_\circ$ 在一个闭\Term{geometric-point}{几何点}处的\Term{strict-localization}{严格局部化}。以 $x_\circ$ 的一个\Term{neighborhood}{邻域}代替 $X_\circ$ 后，可以取 $X_\circ$ \Term{affine}{仿射}、\Term{dimension}{维数} $\leq d$。于是 $X$ 是 $X_\circ$ 上\Term{affine}{仿射}\Term{etale}{平展}\Term{scheme}{概形} $X_i$ 的\Term{projective-limit}{逆极限}（第 VIII 讲 4.5）。

设 $F$ 是 $U$ 上的\Term{torsion-sheaf}{挠层}。须证明 $q>d$ 时 $H^q(U,F)=0$，同样只须处理 $F$ \Term{constructible}{可构造}的情形。此时数据 $f$、$F$、$U$ 来自某个 $X_i$ 上的类似数据（第 IX 讲 2.7）。得到 $U_i=X_i-V(f_i)$ 上的\Term{sheaf}{层} $F_i$，在 $U$ 上“诱导”$F$。由于 $U=\varprojlim U_i$，由第 VII 讲 5.8，有
\[
H^q(U,F)=\varinjlim H^q(U_i,F_i).
\]
因此只须证明如下事实：对任意 $i$，存在\Term{commutative-diagram}{交换图}
\[
\begin{tikzcd}
U\arrow[rr]\arrow[dr]&&U''\arrow[dl]\\
&U_i&
\end{tikzcd}
\]
其中 $U\to U_i$ 是给定\Term{arrow}{箭头}，且 $\operatorname{cd}(U'')<d$。事实上，这将意味着，对 $q>d$，$H^q(U_i,F_i)$ 的每个元素在 $H^q(U,F)$ 中的像都为零。

假定 $i=0$，考虑由 $f_\circ\in\Gamma(X_\circ,\mathcal O_{X_\circ})$ 给出的\Term{morphism}{态射} $g:X_\circ\to\mathbf E_k^1=Y_\circ$。设 $Y''$ 是 $\mathbf E_k^1$ 在原点上方某个\Term{geometric-point}{几何点}处的\Term{strict-localization}{严格局部化}，因而 $Y''$ 是\Term{discrete-valuation-ring}{离散赋值环}的谱；置 $X''=X_\circ\times_{Y_\circ}Y''$。
\SourcePage{118}
由于 $X$ \Term{strictly-local}{严格局部}，\Term{morphism}{态射} $X\to X_\circ$ 经由 $X''$ \Term{resolution}{分解}。而且，\Term{open-subset}{开子集} $U''=X''\times_{X_\circ}U_\circ$ 正是 $X''/Y''$ 的\Term{generic-fiber}{泛纤维}。因此它是 $Y'$ 的\Term{generic-point}{泛点}上的\Term{affine-algebra}{仿射代数}\Term{scheme}{概形}，显然 $\dim U''\leq d-1$。

验证 $\operatorname{cd}U''\leq d$（这就完成证明）。设 $K$ 是 $Y''$ 在\Term{generic-point}{泛点}处的\Term{residue-field}{剩余域}。它是\Term{cohomological-dimension}{上同调维数}为 $1$ 的域，即 $\operatorname{cd}(G)=1$，其中 $G$ 是 $K$ 的\Term{separable-closure}{可分闭包} $\overline K$ 的 Galois 群（第 X 讲 2.2）。设
\[
\overline U''=U''\times_{\operatorname{Spec}K}\operatorname{Spec}\overline K.
\]
由对 $d$ 的归纳假设，$\operatorname{cd}\overline U''\leq d-1$（对于\Term{dimension}{维数} $\leq d$ 的\Term{scheme}{概形}，这是第 X 讲 2.2）。于是 $\operatorname{cd}U''\leq d$ 由 Hochschild--Serre \Term{spectral-sequence}{谱序列}（第 VIII 讲 8.4）
\[
H^q(G,H^{q'}(\overline U,\overline F))\Longrightarrow H^n(U,F)
\]
得到，故引理 4.5 成立。
\end{proof}
\SourcePage{119}
% 原 PDF 第 119 页为空白页。
